\documentclass[11pt]{article}
\usepackage[margin=1in]{geometry}
\usepackage{amsmath,amssymb,amsthm}
\usepackage{hyperref}

\newtheorem{theorem}{Theorem}
\newtheorem{remark}[theorem]{Remark}

\title{Bounded Almost-Sure Noise for Scalar Thresholding}
\author{Research Architecture Runtime \\ \texttt{feature/lean-gated-operator-papers}}
\date{2026-07-25}

\begin{document}
\maketitle

\begin{abstract}
For \(\widetilde A_T(x)=\mathbf{1}\{x+\xi\ge T\}\) with
\(\lvert\xi\rvert\le\eta\) almost surely, pathwise application of the
deterministic threshold-preservation theorem yields almost-sure output
invariance outside the unstable band \([T-\eta,T+\eta)\).
We explicitly \textbf{do not} claim the full Sparse Vector privacy theorem.
Publication is gated on Lean \texttt{LEAN\_FULL} for
\texttt{Research.Operators.Threshold.BoundedNoise}
(pathwise surrogate of a.s.\ bounded noise;
\texttt{REAL\_MATHLIB}).
\end{abstract}

\section{Operator}
Same equality convention \(x\ge T\) passes.
Noise is query-side only in the verified claim; threshold noise and
\(W=\xi-\zeta\) appear as portfolio candidates.

\section{Main theorem}
\begin{theorem}[Bounded-noise threshold]\label{thm:noisy}
Let \(\eta\ge 0\) and \(\lvert\xi\rvert\le\eta\) a.s. Then:
\begin{enumerate}
\item if \(x\ge T+\eta\) then \(\widetilde A_T(x)=1\) a.s.;
\item if \(x<T-\eta\) then \(\widetilde A_T(x)=0\) a.s.;
\item if \(x\in[T-\eta,T+\eta)\) and \(\eta>0\), the random output need not be a.s.\ constant:
under the two-point law \(\xi\in\{-\eta,+\eta\}\) with positive mass on each atom,
both outputs occur with positive probability on this band.
(At the endpoint \(x=T-\eta\), the \(+\eta\) atom yields \(x+\eta=T\), which passes;
at interior points both \(x\pm\eta\) land strictly on opposite sides of \(T\).
If \(\eta=0\) the unstable band is empty and part~(3) is vacuous.)
\end{enumerate}
\end{theorem}

\begin{proof}
Parts (1)--(2): apply the deterministic theorem pathwise with \(\varepsilon:=\eta\).
Part (3): for the two-point law on \(\{\pm\eta\}\), the images \(x\pm\eta\) realize
both sides of the cut on the unstable band (same adversaries as deterministic
sharpness, including the equality pass at \(x=T-\eta\)).
\end{proof}

\begin{theorem}[Noisy sharpness]\label{thm:noisy-sharp}
There exist admissible laws of \(\xi\) with \(\lvert\xi\rvert\le\eta\) a.s.\ such that
if \(x\in[T,T+\eta)\) then \(\mathbb{P}(\widetilde A_T=0)>0\), and if
\(x\in[T-\eta,T)\) then \(\mathbb{P}(\widetilde A_T=1)>0\).
At \(x=T+\eta\), \(\widetilde A_T=1\) a.s.\ for every such \(\xi\).
\end{theorem}

\section{Laplace utility identities (not a DP proof)}
For \(\xi\sim\mathrm{Lap}(0,b)\) with margin \(m=x-T\),
\[
\mathbb{P}(x+\xi\ge T)=
\begin{cases}
1-\tfrac12 e^{-m/b}, & m\ge 0,\\
\tfrac12 e^{m/b}, & m<0.
\end{cases}
\]
The repository checks this closed form at \(m\in\{0,\pm\ln 2\}\) (with \(b=1\))
and margin monotonicity.
This is \textbf{not} a differential-privacy accounting proof.
Feedback limitation: \texttt{LAPLACE\_CDF\_IDENTITY\_NOT\_DP\_PROOF}.

\section{Limitations}
Bounded noise \(\neq\) Sparse Vector; no adaptive positive-release bound.
Lean proof is a \textbf{pathwise} surrogate of a.s.\ bounded noise
(\texttt{REAL\_MATHLIB}); Laplace identities remain engineering checks,
not a DP proof (\texttt{LAPLACE\_CDF\_IDENTITY\_NOT\_DP\_PROOF}).

\section{Verifier summary}
Publication requires \texttt{LEAN\_FULL} for
\texttt{Research.Operators.Threshold.BoundedNoise}
(\texttt{bounded\_noise\_preservation}, \texttt{bounded\_noise\_sharpness}).
Certificate:
\texttt{lean/certificates/thresholding/bounded-noise-threshold/}.
Prior computational-only paper archived under \texttt{\_archive/v0-computational/}.

\section*{References}
\begin{thebibliography}{9}
\bibitem{dwork2014}
Dwork and Roth, Algorithmic Foundations of Differential Privacy, 2014.
\bibitem{repo}
\texttt{implementation/src/operators/thresholding/}.
\end{thebibliography}

\end{document}
